\documentclass{article}

\usepackage[preprint]{neurips_2026}

\usepackage[utf8]{inputenc}
\usepackage[T1]{fontenc}
\usepackage{hyperref}
\usepackage{url}
\usepackage{booktabs}
\usepackage{amsfonts}
\usepackage{amsmath}
\usepackage{nicefrac}
\usepackage{microtype}
\usepackage{xcolor}

\title{Inference Pipeline for Improving Mathematical Reasoning on Qwen3-4B-Thinking-2507}

\author{%
  Fong-Yu Lin \\
  UC San Diego \\
  \texttt{[email]} \\
  \And
  Anthony Nguyen \\
  UC San Diego \\
  \texttt{[email]} \\
}

\begin{document}

\maketitle

\begin{abstract}
We describe our CSE 151B competition submission pipeline for mathematical reasoning with \texttt{Qwen/Qwen3-4B-Thinking-2507}. The final submitted system runs on H200 hardware using pure model inference, thinking mode, self-consistency with $K=5$, adaptive retry, and two private shards. This report documents the final H200 K=5 split run, the engineering decisions behind it, the validation results, and the QLoRA experiments we rejected.
\end{abstract}

\section{Introduction}

The competition task is to generate mathematical solutions for a private test set and submit a CSV containing an \texttt{id} and a model-generated \texttt{response} for each problem. A useful response must not only contain correct reasoning, but must also expose the final answer in a format that can be extracted reliably, typically using \verb|\boxed{...}|. The private set contains 943 examples, and private labels are not available, so final validation must focus on structural correctness: complete row coverage, no duplicate IDs, no empty responses, and no manual modification of model answers.

Our project explored several possible directions, including prompt engineering, self-consistency, Program-of-Thought, QLoRA adapters, and public validation sweeps. The final submitted run is the H200 second round with $K=5$ and adaptive retry, split into two shards covering indices 0--472 and 472--943. The two shards merged cleanly into a complete 943-row CSV, so this K=5 run is the submitted system.

\section{Submitted system}

\subsection{Model and hardware}

The submitted system uses the required base model, \texttt{Qwen/Qwen3-4B-Thinking-2507}. No fine-tuned checkpoint, external model, external API, calculator, or generated-code execution is used at private inference time. The H200 run uses the legal pure-inference pipeline exposed through \texttt{run\_inference.py}.

\begin{table}[h]
\centering
\caption{Earlier A30 backup notebook used during development, not submitted.}
\label{tab:k1-config}
\begin{tabular}{ll}
\toprule
Setting & Value \\
\midrule
Notebook & \texttt{inference\_sc\_k1\_private (1).ipynb} \\
Model & \texttt{Qwen/Qwen3-4B-Thinking-2507} \\
Data path & \texttt{data/private.jsonl} \\
Private examples & 943 \\
GPU & NVIDIA A30 \\
Precision & bfloat16 \\
Samples per question & $K=1$ \\
Maximum generation tokens & 16,384 \\
Maximum model length & 32,768 \\
Temperature / top-p / top-k & 0.6 / 0.95 / 20 \\
Chunked prefill & enabled \\
\bottomrule
\end{tabular}
\end{table}

\begin{table}[h]
\centering
\caption{Configuration of the final submitted H200 K=5 private run.}
\label{tab:h200-config}
\begin{tabular}{ll}
\toprule
Setting & Value \\
\midrule
Model & \texttt{Qwen/Qwen3-4B-Thinking-2507} \\
Data path & \texttt{data/private.jsonl} \\
Shard 1 & indices 0--472, \texttt{submission\_part\_000\_472.csv} \\
Shard 2 & indices 472--943, \texttt{submission\_part\_472\_943.csv} \\
Samples per question & $K=5$ \\
Maximum generation tokens & 24,576 \\
Maximum model length & 32,768 \\
GPU memory utilization & 0.90 \\
Maximum sequences & 32 \\
Maximum batched tokens & 32,768 \\
Generation chunk size & 64 \\
Prefix caching & disabled \\
Adaptive retry & enabled, \texttt{retry\_k=2}, \texttt{retry\_max\_tokens=32768} \\
\bottomrule
\end{tabular}
\end{table}

\subsection{Prompting}

We use two task-specific prompt templates. For free-form problems, the system prompt asks the model to give one answer per \texttt{[ANS]} placeholder, in order, inside a single \verb|\boxed{}| expression, without units or prose inside the box. For multiple-choice problems, the system prompt asks for exactly one capital letter inside \verb|\boxed{}| and explicitly forbids putting formulas, numeric values, or option text inside the final box.

The prompt design is intentionally simple. We found that broad instructions such as ``solve carefully'' are not sufficient by themselves: the final answer format is an important part of the task. The prompts therefore emphasize the final answer contract while still allowing the Qwen thinking model to produce a full reasoning trace.

\subsection{Inference flow}

The inference flow is checkpointed in the H200 split run:

\begin{enumerate}
  \item Load all private examples from \texttt{data/private.jsonl}.
  \item Split the 943 examples into recoverable chunks or index shards.
  \item Select assigned sections, chunks, or index ranges for the current runner.
  \item Render each problem using the Qwen chat template with thinking mode enabled.
  \item Generate $K=5$ responses per problem with vLLM.
  \item Extract the final nested \verb|\boxed{...}| answer for diagnostics.
  \item Write each completed chunk immediately to a CSV checkpoint.
  \item Merge all chunk CSVs into the final Kaggle submission.
\end{enumerate}

The chunked design is important for deadline reliability. A full private run can be interrupted by wall-clock limits, process crashes, or GPU availability. By saving raw JSONL checkpoints after each generation chunk, the team can rerun the same shard command and resume completed questions.

\section{Validation}

Because private labels are unavailable, private accuracy cannot be measured locally. We therefore validate the submission structurally. The final merged CSV contains exactly 943 rows, one for each private ID, with no missing IDs, no extra IDs, no duplicates, and no empty responses. The repository includes \texttt{merge\_submission\_shards.py}, which performs these checks before writing the final CSV. The H200 merge inputs were \texttt{results/submission\_part\_000\_472.csv} and \texttt{results/submission\_part\_472\_943.csv}.

\begin{table}[h]
\centering
\caption{Final K=5 submission validation checklist.}
\label{tab:validation}
\begin{tabular}{lll}
\toprule
Check & Expected & Final value \\
\midrule
Rows in final CSV & 943 & 943 \\
Missing private IDs & 0 & 0 \\
Extra IDs & 0 & 0 \\
Duplicate IDs & 0 & 0 \\
Empty responses & 0 & 0 \\
Kaggle/private score & unavailable locally & unavailable locally \\
\bottomrule
\end{tabular}
\end{table}

\section{Stronger pipeline design}

The submitted pipeline in \texttt{run\_inference.py} uses the base model and legal private-inference constraints, while adding self-consistency and adaptive retry. This is the pipeline used for the H200 K=5 final run.

The submitted settings are:

\begin{itemize}
  \item $K=5$ independent samples per question.
  \item \texttt{max\_tokens=24576}.
  \item \texttt{max\_model\_len=32768}.
  \item \texttt{temperature=0.6}, \texttt{top\_p=0.95}, \texttt{top\_k=20}.
  \item adaptive retry when no boxed answer is found, when votes tie, or when a sample is length-truncated.
  \item raw JSONL checkpoints and metadata summaries.
  \item public scoring when the input file includes public labels.
\end{itemize}

This method should be more accurate in principle because majority voting reduces single-sample errors and retry gives low-confidence questions additional chances. The cost is substantially higher wall-clock time and memory pressure than a single-sample run, which is why we used shard-level checkpointing and merge validation.

\section{Public sweep result}

To choose the final configuration objectively, we used \texttt{sweep\_inference\_configs.py}. It evaluates several public-set candidates and writes a summary table.

\begin{table}[h]
\centering
\caption{Public validation sweep configurations.}
\label{tab:sweep}
\begin{tabular}{lll}
\toprule
Name & Main setting & Purpose \\
\midrule
A & $K=5$, \texttt{max\_tokens=24576} & Selected final setting \\
B & $K=7$, \texttt{max\_tokens=24576} & More voting at higher cost \\
C & $K=3$, \texttt{max\_tokens=24576} & Lower-cost self-consistency \\
D & $K=5$, \texttt{max\_tokens=16384} & Lower reasoning budget \\
E & $K=5$, \texttt{retry\_k=4} & More retry samples for failures \\
\bottomrule
\end{tabular}
\end{table}

The selected public configuration scored 103/150 overall on the public 0--150 sweep, with 47/56 MCQ accuracy and 56/94 free-form accuracy. It also had high boxed-answer and thinking-completion coverage, so we used it for the final private run.

\section{Exploratory work not used for submission}

We explored Program-of-Thought prompting, where the model writes Python code and the system executes it to obtain an answer. This can improve arithmetic reliability in local analysis, especially for computational problems. However, it executes generated code through a subprocess and is therefore not part of our private submission path.

\subsection{QLoRA experiments on \texttt{yang-test}}

We also explored QLoRA adapter training on the \texttt{yang-test} branch. This work produced useful infrastructure: a scriptable QLoRA SFT pipeline, public train/dev splitting, NuminaMath-CoT streaming support, Transformers-based adapter evaluation, raw-result scoring, and a Markdown result tracker. However, the adapter results did not justify using QLoRA for the final private submission.

\begin{table}[h]
\centering
\caption{Selected QLoRA results from \texttt{yang-test}. These results led us to reject the adapter path for the submitted system.}
\label{tab:qlora-results}
\begin{tabular}{llll}
\toprule
Experiment & Base control & Adapter result & Decision \\
\midrule
Public smoke adapter, $n=10$ & 4/10 & 3/10 & Reject \\
Numina 5k, 2048-token eval, $n=50$ & 17/50 & 20/50 & Inconclusive \\
Numina 5k, 4096-token eval, $n=50$ & 26/50 & 17/50 & Reject \\
\bottomrule
\end{tabular}
\end{table}

The public-smoke adapter was useful as an engineering test: it verified that 4-bit loading, LoRA attachment, training, adapter saving, and evaluation all worked. But it performed worse than the base control on the held-out smoke split. Qualitatively, it also tended to produce much shorter answers, suggesting that answer-only public supervision taught output brevity more than robust reasoning.

The NuminaMath-CoT adapter was a more serious attempt because Numina provides worked solutions rather than only final answers. At a 2048-token evaluation budget it had a small overall gain, 20/50 versus 17/50, but the gain was not stable: free-form accuracy was worse than the base control. When we increased the evaluation budget to 4096 tokens under matched Transformers settings, the base model reached 26/50 while the Numina 5k adapter reached only 17/50. This was the strongest evidence that the adapter should not be used for final submission.

There were also engineering reasons to avoid QLoRA under deadline pressure. In our setup, vLLM was useful for fast base-model prompt experiments but could not run the Qwen3 LoRA adapter path reliably, so fair adapter comparisons required the slower Transformers backend. Since our final private run needed to finish reliably, the base-model inference path was safer.

\section{Discussion}

The main strength of our submission strategy is reliability. The H200 K=5 run keeps a legal inference-only design while adding stronger sampling, voting, and retry. Shard-level checkpointing makes the run recoverable, and structural validation ensures the final CSV is well formed.

The main weakness is that private accuracy cannot be measured locally, so public validation is the only available proxy for choosing between configurations. The K=5 approach also requires more compute time and a careful shard merge. Finally, our QLoRA results show that fine-tuning is not automatically beneficial: an adapter can reduce the model's useful reasoning behavior if the training data or token budget is not well matched to the evaluation setting.

\section{Future work}

Our immediate future work is operational: submit the validated H200 K=5 CSV and keep the repository aligned with that exact configuration. If more time were available, we would consider an even stronger private run with $K=7$ or a larger retry budget, but only after public validation and structural checks.

Longer term, the most promising direction is self-distillation from correct reasoning traces. A strong base model can generate multiple public-train responses; correct and well-formatted traces can then be filtered and used for supervised fine-tuning. This is likely more useful than training only on final answers, because it teaches the model the reasoning process rather than only the output format. If we revisit QLoRA, we should train on those filtered correct traces and only use the adapter if it beats the base model under the same split, backend, prompt mode, token budget, and sampling settings.

\section{Conclusion}

Our final strategy is deadline-aware: submit the validated H200 K=5 split run. This gives us a stronger completed submission, a reproducible inference path, and a clear record of the validation process.

\end{document}
